\section{Reformular el cáncer como un problema de modelo del mundo}

Las dos secciones anteriores siguen siendo válidas para robots, video o voz. A partir de esta sección se explica por qué la biología tumoral no es una aplicación elegida al azar: es un sistema dinámico multiescala y, a la vez, dispone de varios sensores complementarios pero costosos. Lo que de verdad importa en la clínica no es «a qué clase pertenece esta imagen», sino «si a este paciente se le administra cierto tratamiento, ¿el tumor cambiará de forma favorable?».

\subsection{Correspondencia entre el mundo macroscópico y el mundo tumoral}

Del paisaje urbano al tumor, la estructura se mantiene: el estado real no puede observarse por completo, los sensores aportan evidencia parcial, las acciones modifican el futuro y la consulta determina el objetivo de la simulación. La diferencia es que las variables de estado del tumor abarcan varias escalas: paciente, órgano, tejido, célula y molécula. A continuación, mediante tres figuras (el mundo macroscópico, el microambiente inmunitario y el simulador de pacientes), los componentes abstractos se asignan paso a paso a objetos medibles y a preguntas clínicas.

\lecturefigure{slide-092-macroscopic-world-model.jpg}{Modelo del mundo macroscópico: cámara, micrófono y texto responden en conjunto a consultas sobre acciones}{00:22:46}

Esta diapositiva de repaso recuerda que el valor de un world model proviene del ciclo cerrado: los sensores no sirven para reunir la mayor cantidad posible de datos, sino para responder consultas relacionadas con las acciones. Si una modalidad es costosa pero no cambia el orden de preferencia entre tratamientos, su valor marginal puede ser muy bajo; si una medición poco frecuente permite distinguir dos subpoblaciones con respuestas al tratamiento totalmente opuestas, resulta fundamental.

\lecturefigure{slide-094-cancer-immunology.jpg}{Inmunología del cáncer en una diapositiva: el tumor y el sistema inmunitario interactúan de forma continua en el microambiente}{00:25:04}

Las células de colores de la figura no son decorativas. El Tumor microenvironment (microambiente tumoral) incluye células tumorales, linfocitos T, linfocitos B, macrófagos y estroma, entre otros; el sistema inmunitario puede reconocer y eliminar células anómalas, pero también puede ser suprimido, excluido o reprogramado por el tumor. Por ello, la respuesta al tratamiento depende de la combinación de tipos celulares, posición espacial, estados y vías de señalización.

\begin{knowledgebox}{Primera aparición: inmunoterapia tumoral}
La Cancer immunotherapy (inmunoterapia contra el cáncer) no busca destruir directamente todas las células tumorales, sino eliminar la supresión inmunitaria, potenciar el reconocimiento o modificar la función de las células inmunitarias. Un mismo fármaco produce respuestas muy distintas en pacientes diferentes, y eso es justamente lo que un modelo del mundo condicionado al paciente pretende explicar.
\end{knowledgebox}

\lecturefigure{slide-095-patient-simulator.jpg}{Objetivo ideal: partir del tejido del paciente para simular la respuesta a nivel de tejido después del tratamiento}{00:25:20}

En la figura, el tejido del paciente se introduce en el world model y luego se pregunta «¿remitirá el tumor después de administrar el fármaco?». Esto puede formalizarse como
\[
p_\theta(R\mid X_{\text{patient}},a_{\text{drug}}),
\]
donde $R$ es la respuesta, $X_{\text{patient}}$ es el estado multimodal del paciente y $a_{\text{drug}}$ es la intervención farmacológica. La fórmula parece sencilla; lo realmente difícil es que $X$ no puede medirse por completo, que el efecto del fármaco tarda mucho tiempo en manifestarse, que los datos de entrenamiento carecen de intervenciones aleatorizadas y que las etiquetas de respuesta clínica son escasas.

\lecturefigure{slide-096-dataset-scale.jpg}{Escala de los datos en el momento de la clase: pacientes, cortes, píxeles y mediciones espaciales crecen en conjunto}{00:25:21}

Esta diapositiva de cifras debe leerse en su contexto temporal: describe una instantánea de la investigación en el momento de la clase, en mayo de 2025, no la situación de 2026. Más importante que los conteos concretos es la diferencia entre niveles: del orden de miles de pacientes, miles de cortes, una enorme cantidad de píxeles y de puntos de medición espacial implica que el «tamaño de muestra» no puede evaluarse solo por el número de pacientes ni solo por el número de tokens. La independencia entre pacientes determina la generalización; las células y los píxeles determinan la precisión estadística local.

\teachervoice{00:22:46--00:25:23, el docente condensa el problema en una sola consulta clínica: si a este paciente se le administra cierto tratamiento, ¿remitirá el tumor? Esto restringe los datos y la validación mejor que «construir un modelo fundacional para el cáncer».}

\begin{warningbox}{Un modelo del mundo no equivale a un sistema de decisión clínica}
Aunque el modelo pueda predecir cambios en el tejido, aún deben considerarse la toxicidad, la dosis, la farmacocinética, las comorbilidades del paciente y la incertidumbre. La simulación a nivel de tejido es evidencia candidata, no una prescripción directa. En estas notas, el «descubrimiento de fármacos» debe entenderse como una etapa del proceso de descubrir y jerarquizar hipótesis.
\end{warningbox}

\subsection{Resumen de la sección}

La biología tumoral puede reformularse como un sistema dinámico parcialmente observable, sujeto a intervención y multiescala. El objetivo real es la respuesta al tratamiento condicionada al paciente, no una clasificación estática. Este objetivo exige que el modelo maneje a la vez modalidades complementarias, estructura espacial, un número limitado de pacientes y una necesidad muy alta de validación causal.
